\sscp{Lexical features in Austronesian and Papuan languages}{15-04-01}
Throughout this study we have encountered several words that
are similar in both form
and meaning in the Wallacean region
and are found in both Austronesian and Papuan languages.

%In \srf{02-03-01-07}, we discussed several important food crops
%that originated in the New Guinea area
%and spread westward before the arrival of people speaking Austronesian languages.
%We specifically looked at the pseudo-etymon of {\#}muku ‘banana’.
%We saw that terms relating to this form were found in several unrelated Papuan families.
%and we saw that terms relating to this form were found in many
%Austronesian languages in several Wallacean subgroups extending
%as far west as western Flores (see \frf{02-20} and appendix \ref{ch:BanTer}).
%We noted that the distribution of terms relating to this form
%was limited to the southern string of islands from the Bird’s
%Neck of Papua to Flores
%and Sumba, and not found across central Maluku except close to the Bird’s Neck.
%We noted that reflexes are found in some nodes of \tsc{Seram-Tanimbar-Bomberai}
%close to the Bird’s Neck
%and not in others.
%Reflexes are found in some nodes of \tsc{Timor-Babar}
%and not in others.
%\citet[478,486]{UsherSchapper2022} more specifically assert that the words in Austronesian
%languages that we associate with {\#}muku ‘banana’ ultimately
%link back to the Papuan Proto-Greater West Bomberai *maŋgʷVl ‘banana’.
%Such Papuan loans in Austronesian languages are clear evidence of contact.
%The greater diversity of forms in Papuan languages
%and the greater consistency of forms in Austronesian languages
%suggests that the word was borrowed only once or twice into Austronesian
%languages and then spread between them by local contact or diffusion.
%We noted that reflexes of this word coexist side-by-side with
%neighbouring languages that retain PMP *punti ‘banana’,
%often in the same nodes of the same subgroups.
%and there are several other etyma also meaning ‘banana’ found in the region.
%The spread of {\#}muku ‘banana’ into some Wallacean subgroups
%but not others is evidence that new words can
%and do spread through local links within this region of long-standing contact and trade.
%The distribution of this word across parts of the southern islands,
%but not across central Maluku shows that this word cannot have
%been borrowed into and then spread from a common parent language (such as PCMP).
%and the presence of reflexes of this word in some nodes of certain
%Wallacean subgroups, but not in other nodes of the same subgroup shows that {\#}muku
%‘banana’ cannot be assigned to the proto-language of those individual
%subgroups and is further evidence of contact.
%Each of these issues is significant.
%and the patterns associated with {\#}muku ‘banana’ provide a
%model for possible alternate explanations for the dispersal (or
%diffusion) of other problematic words in the region that have
%been assigned to PCMP and PCEMP.

%In \srf{14-03} we showed that there is a chain of incrementally
%resemblant forms meaning `cuscus' stretching from Sulawesi through north Maluku
%into the Bird’s Neck region of Papua (see \frf{14-03}).
%The form **kusay ‘bear cuscus’ is found in several Austronesian
%subgroups throughout much of Sulawesi; **kusoy ‘bear cuscus’
%is found in some Austronesian languages in the north fairly close
%to Halmahera; proto-North Halmahera *kusor ‘cuscus’ accounts
%for at l\ili{east} seven Papuan languages of north Halmahera; {\#}kisor(o)
%‘cuscus, marsupial’ accounts for some Austronesian languages
%and languages in several Papuan families in the Bird’s Head
%and Bird’s Neck region of Papua;
%and **kindoR(a) or **kVndoR(a) ‘cuscus’ is reflected in some
%languages of the \tsc{Eastern Islands} node of STB
%and perhaps in a few SHWNG languages.
%The close similarities of form
%and meaning across unrelated language families
%and subgroups must be seen as evidence of contact.
%The distribution of *kVndoR(a) ‘cuscus’ limited to only one node
%of STB and not found other nodes of STB nor in any other Wallacean subgroup
%means this etymon cannot be assigned to STB,
%nor can it be assigned to the purported higher level parent languages
%of PCMP or PCEMP.
%
%In \srf{14-04} we showed that **mantəR ‘cuscus’
%and incrementally resemblant forms (such as **məntəR,
%**mədaR) are found in many Austronesian languages stretching
%from \tsc{\ili{Muna}-Buton} in Sulawesi throughout much of linguistic
%Wallacea, and into Papua, where it is also found in several widely scattered
%Papuan languages and sometimes with a gloss shifted there to other marsupials
%such as the tree kangaroo (see \frf{14-05}).
%Reflexes of **mantəR ‘cuscus’ are not found in \tsc{\ili{Bima}-Lembata},
%which is outside the range of marsupials.
%The term has no reflexes in the \tsc{Eastern Islands} node of
%STB which instead reflects **kVndoR(a) ‘cuscus’.
%
%In \srf{14-04-04}, following \cite[269]{Schapper2011},
%we noted the similarities in form between Austronesian
%**mantəR ‘cuscus’ and neighbouring Papuan
%Proto-Greater West Bomberai *mandəl ‘flying
%fox’, both referring to edible animals hunted in the treetops.
%Those two groups are intermixed in the eastern parts of our target region.
%Within STB, \tsc{East Seram} languages explicitly link cuscus with bats.
%Far to the west,
%Proto-Timor-\ili{Alor}-\ili{Pantar} *madəl ‘flying fox’,\footnote{
		%\tsc{Timor-\ili{Alor}-Pantar}
		%is proposed to be a branch of \tsc{Greater West Bomberai} \citep{UsherSchapper2022}.}
 %is fairly resemblant of nearby Proto-\ili{Central} Timor **mantəR
%‘cuscus’, or Proto-Timor-\ili{Babar} **mantəR/**məntəR ‘cuscus’.
%\citet[486]{UsherSchapper2022} also see the Papuan forms as the
%source of loanwords into several groups of Austronesian languages
%across Wallacean meaning ‘bat’.

Among terms for `cuscus', in \srf{14-03} we showed
that there is a chain of incrementally \it{k}-initial
resemblant words stretching from Sulawesi through north Maluku
into the Bird’s Neck region of Papua (see \frf{14-03}).

Likewise, in \srf{14-04} we showed that **mantəR ‘cuscus’
and incrementally resemblant forms (such as **məntəR,
**mədaR) are found in many Austronesian languages stretching
from \tsc{\ili{Muna}-Buton} in Sulawesi throughout much of linguistic
Wallacea, and into Papua, where it is also found in several widely scattered
Papuan languages and sometimes with a gloss shifted there to other marsupials
such as the tree kangaroo (see \frf{14-05}).
In \srf{14-04-04}, following \cite[269]{Schapper2011},
we noted the similarities in form between Austronesian
**mantəR ‘cuscus’ and neighbouring Papuan
Proto-Greater West Bomberai *mandəl ‘flying fox’,
both referring to edible animals hunted in the treetops.
\citet[486]{UsherSchapper2022} also see the Papuan forms as the
source of loanwords into several groups of Austronesian languages
across Wallacea meaning ‘bat’.

In \srf{01-08} and \srf{02-01} we noted in passing the
presence of {\#}muku `banana' in both Austronesian and
Papuan languages of Wallacea.

So consequently we have clear evidence of new words appearing
in the Austronesian languages in the target region,
some of which are from Papuan sources,
which spread unevenly across several Wallacean subgroups through
contact, are interspersed with other terms,
and which cannot be assigned to a common Austronesian parent
language such as PMP, PCMP or PCEMP.

In the following sections we explore in
detail several words and phrases shared
between Austronesian and Papuan languages in our target region.